\chapter{Percutaneous Nephrolithotomy (PCNL)}

\section*{Introduction \& Clinical Indications}

Percutaneous nephrolithotomy (PCNL) is a minimally invasive surgical procedure utilized for the removal of large renal calculi (kidney stones) through a small incision in the skin. This technique is particularly advantageous for stones that are greater than 2 cm in size, those located in complex anatomical positions, or when other less invasive methods, such as extracorporeal shock wave lithotripsy (ESWL) or ureteroscopy, have failed. 

Historically, the development of PCNL can be traced back to the 1970s, with significant advancements in imaging and instrumentation leading to its widespread adoption. The pathophysiological rationale for PCNL lies in the need to effectively manage obstructive uropathy and prevent complications such as infection, renal damage, and chronic pain associated with large stones.

Clinical indications for PCNL include:

\begin{itemize}
    \item Large renal stones (>2 cm).
    \item Staghorn calculi.
    \item Obstructive uropathy due to stones.
    \item Failed previous treatments (e.g., ESWL).
    \item Anatomical abnormalities complicating stone removal.
    \item Patients with a solitary kidney.
    \item Recurrent stone formation despite conservative management.
\end{itemize}

PCNL is often considered a first-line treatment for large stones, particularly in patients with significant comorbidities or those who prefer a minimally invasive approach. In cases where PCNL is not feasible, alternative options may include open surgery or ureteroscopy.

\section*{Relevant Surgical Anatomy}

A thorough understanding of the relevant surgical anatomy is crucial for the successful execution of PCNL. The kidneys are retroperitoneal organs located in the upper abdomen, with the right kidney typically positioned slightly lower than the left due to the presence of the liver.

Key anatomical structures include:

\begin{itemize}
  \item \textbf{Collecting system}: The renal pelvis and calyces collect urine and provide the target for safe access.
  \item \textbf{Vascular anatomy}: Segmental renal arteries are functionally end arteries; the access tract should avoid major vascular branches where possible.
  \item \textbf{Adjacent structures}: The pleura, colon, spleen, liver, and bowel may be at risk depending on side, stone location, and puncture tract.
  \item \textbf{Surface landmarks}: The 11th or 12th rib, posterior axillary line, and selected posterior calyx guide image-assisted access; the tract is planned from preoperative imaging.
\end{itemize}

Surgical landmarks include the 12th rib, which serves as a guide for the initial incision, and the posterior axillary line, which helps in determining the appropriate access point for nephrostomy.

\section*{Preoperative Workup \& Preparation}

A comprehensive preoperative workup is essential to optimize patient outcomes and minimize complications. Key components include:

\begin{itemize}
  \item \textbf{Testing}: Urinalysis and urine culture, blood count, renal function, and coagulation testing are obtained as clinically indicated. A positive culture requires treatment and procedure planning.
  \item \textbf{Imaging}: Noncontrast CT defines stone burden and surrounding anatomy; ultrasound or MRI-based approaches may be selected when radiation should be avoided.
  \item \textbf{Preparation}: Anticoagulants, antiplatelets, antibiotics, anesthesia risk, and the risks, benefits, alternatives, and possible staged procedures are reviewed. Fasting follows the anesthesia service.
\end{itemize}

Uncorrectable coagulopathy, an unsafe access route, or untreated urinary infection may require postponement or an alternative approach. Obesity, abnormal renal anatomy, solitary kidney, and prior surgery increase technical complexity but are not automatic contraindications.

\section*{Surgical Technique \& Procedure}

The surgical procedure for PCNL involves several critical steps:

\begin{enumerate}
  \item \textbf{Positioning and anesthesia}: The patient is positioned prone or supine according to anatomy, expertise, and airway considerations, and anesthesia is provided.
  \item \textbf{Image-guided access}: A selected calyx is punctured under fluoroscopic or ultrasound guidance and a guidewire is advanced into the collecting system.
  \item \textbf{Tract and nephroscopy}: The tract is dilated and a nephroscope is introduced. Laser, ultrasonic, or pneumatic devices fragment the stone.
  \item \textbf{Extraction and drainage}: Fragments are removed and the collecting system is inspected. A nephrostomy tube, ureteral stent, or tubeless completion is selected according to residual stone, bleeding, infection, and tract findings.
\end{enumerate}

\section*{Complications \& Risks}

As with any surgical procedure, PCNL carries inherent risks. General surgical complications include:

\begin{itemize}
    \item Hemorrhage.
    \item Infection.
    \item Anesthesia-related complications.
\end{itemize}

Procedure-specific risks include:

\begin{itemize}
    \item Injury to surrounding organs (e.g., colon, spleen, liver).
    \item Pleural injury leading to pneumothorax.
    \item Postoperative fever.
    \item Urinary leakage.
    \item Nephrostomy tube complications (e.g., blockage, dislodgement).
    \item Stone recurrence.
\end{itemize}

Understanding these risks is essential for informed consent and patient management.

\section*{Postoperative Care \& Recovery}

Postoperative care begins in the post-anesthesia care unit (PACU), where patients are monitored for vital signs, pain control, and any immediate complications. 

Expected postoperative course includes:

\begin{itemize}
  \item Provide analgesia and monitor urine output, hematuria, fever, sepsis, anemia, renal function, and respiratory symptoms.
  \item Review stone analysis and any renal pelvis or stone culture when obtained; pathology is not routinely relevant to a calculus.
  \item Follow-up imaging and metabolic stone evaluation are individualized according to residual fragments, stone composition, recurrence risk, and renal function.
\end{itemize}

\section*{Outcomes \& Prognosis}

Stone-free outcomes vary with stone burden, anatomy, number of tracts, and staged procedures. PCNL is recommended as first-line treatment for many renal stones larger than 2\,cm, while ureteroscopy or shock-wave lithotripsy may be alternatives when PCNL is unsuitable.

Recurrence rates vary but are estimated to be around 20\% to 30\% within five years. Epidemiological studies indicate that factors such as metabolic abnormalities and dietary habits play a significant role in stone formation.

Landmark trials, such as the Clinical Research Office of the Endourological Society (CROES) study, have provided valuable insights into the efficacy and safety of PCNL, establishing it as a standard of care for large renal stones.

\section*{Additional Clinical Considerations}
PCNL is selected by stone burden and anatomy, not by size alone. A stone larger than 2 cm, a staghorn or branching stone, a lower-pole stone with unfavorable anatomy, or a stone with high density may be difficult to clear with shock-wave lithotripsy. Flexible ureteroscopy can be an alternative for selected patients, particularly when bleeding risk or access anatomy makes percutaneous treatment less attractive, although it may require multiple stages. The choice should consider renal function, infection risk, body habitus, anticoagulation, patient preference, and local expertise.

Infection prevention is central. A urine culture should be obtained and a positive culture treated before elective access when possible. Infected obstruction requires urgent drainage with a nephrostomy or ureteral stent and antibiotics; definitive stone clearance is generally delayed until sepsis has been controlled. Stones may harbor bacteria even when bladder urine is sterile, so renal pelvic urine and stone cultures can help guide treatment in a febrile or septic patient. Fever, tachycardia, hypotension, rigors, or confusion after PCNL warrants prompt evaluation for systemic infection rather than routine reassurance.

Access planning balances clearance with renal and adjacent-organ safety. The preferred calyx and tract are selected from CT or other imaging, and ultrasound may reduce radiation exposure and help identify the pleura, colon, spleen, and liver. Supracostal access can be necessary for upper-pole or staghorn stones but increases the need to assess the pleura and diaphragm. Miniaturized instruments may reduce bleeding in selected cases, but the best tract size and number depend on stone burden, equipment, and the surgeon's plan. A staged procedure is acceptable when attempting complete clearance in one session would increase risk.

Residual fragments are documented rather than assumed to be clinically irrelevant. Follow-up imaging is chosen to balance sensitivity, radiation, and the need to detect obstruction or clinically significant residual stone. Stone analysis and metabolic evaluation can identify calcium, uric acid, infection, cystine, or other stone types. Prevention may include increased fluid intake, dietary sodium reduction, appropriate calcium intake, moderation of animal protein, and targeted medication based on urine studies; advice should not be reduced to a single universal diet.

Bleeding is usually managed conservatively, but a falling hemoglobin, persistent gross hematuria, clot retention, or hemodynamic instability may require transfusion, imaging, angioembolization, or re-intervention. A nephrostomy tube provides drainage and access for a planned second look when residual stone or edema is present. A tubeless or totally tubeless approach is reserved for selected uncomplicated cases with good drainage and hemostasis. The discharge plan should identify the tube or stent, expected hematuria, analgesia, antibiotic plan, and the date or criteria for removal.

\section*{Stone Prevention and Follow-Up}
The procedure treats the current stone burden but does not cure the tendency to form stones. Stone composition, serum studies, and a timed urine collection may identify hypercalciuria, hypocitraturia, hyperoxaluria, low urine volume, uric-acid risk, infection stones, or cystinuria. Prevention is tailored to those findings. Adequate fluid intake is central, while excessive dietary sodium and inappropriate restriction of dietary calcium can increase risk in some calcium-stone formers.

Residual fragments are followed according to size, location, composition, symptoms, and obstruction risk. A small nonobstructing fragment may be observed, whereas a residual staghorn component, infection stone, or obstructing fragment may require staged ureteroscopy, second-look PCNL, or another intervention. Recurrent fever, flank pain, reduced urine output, or persistent hematuria should prompt urgent review. A stent or nephrostomy is removed only after the treating team confirms that drainage is safe.

\section*{Patient Education}
Patients should be told whether a nephrostomy or ureteral stent remains, how it is cared for, and when it will be removed. Some hematuria and flank discomfort may be expected, but clots preventing urine drainage, fever, chills, reduced urine output, severe pain, or breathlessness require urgent review. The discharge plan should include stone analysis, imaging, metabolic evaluation, hydration advice, and the responsible urology contact.

\section*{Documentation and Safety}
The operative record should describe stone burden, side and calyx accessed, imaging guidance, tract size and number, fragmentation method, residual stone, drainage device, blood loss, and complications. The discharge summary should identify culture results, antibiotic instructions, stent or nephrostomy plans, and follow-up imaging. These details are important when a staged procedure or urgent reassessment becomes necessary.

\section*{Practical Follow-Up}
Follow-up confirms drainage, reviews stone and culture results, removes stents or nephrostomy devices at the planned time, and evaluates residual fragments. Fever, rigors, reduced urine output, clot retention, severe flank pain, or breathing difficulty requires urgent urologic assessment.

\section*{Safety Summary}
PCNL is a standard option for many large or complex renal stones, but infection control and safe access come first. Residual fragments, drainage devices, stone cultures, and metabolic prevention determine the follow-up plan.

\section*{Final Safety Note}
PCNL requires infection control, safe image-guided access, and a documented plan for residual stones and drainage devices. Fever, rigors, low urine output, clot retention, severe flank pain, or breathlessness requires urgent urologic assessment.

\section*{Completion Checklist}
Record stone burden, access calyx, tract, residual fragments, cultures, drainage devices, antibiotic plan, and imaging follow-up. Fever, rigors, low urine output, clot retention, severe pain, or breathlessness requires urgent review.

\section*{References}

\begin{enumerate}
    \item Smith, A. D., \& Bultitude, M. F. (2015). \textit{Urology: A Comprehensive Clinical Guide}. New York: Springer.
    
    \item Pearle, M. S., et al. (2016). "Medical Management of Kidney Stones: AUA Guideline." \textit{The Journal of Urology}, 196(4), 1161-1169.
    
    \item Schwartz, S. I., et al. (2019). \textit{Principles of Surgery}. New York: McGraw-Hill.
    
    \item Krambeck, A. E., et al. (2018). "Percutaneous Nephrolithotomy: A Review of the Literature." \textit{Urology}, 112, 1-8.
\end{enumerate}

\clearpage
